\calculussheet{sheet:p18-erotetic-search}{Original erotetic search warrants and imperative procedures}
  {E/Im entries are source warrant templates, not primitive axioms. The eight
  procedure entries reproduce every response case of Figures 5--12, not direct
  question-implication rules. Full compound scenarios retain the unqueried
  \(?\pm|A,B|\). Instances preserve question well-formedness.
  Proofs: \ref{annex:p18.es.search-procedures}.}

\hbrule{rule:p18-es-e-warrant}{E-rule warrant (source schema)}
  {X\qquad E_L(X,Q)}{\text{raise }Q}
  {\(X\) finite; actual evocation gives coverage and individual non-entailment.
   Actual soundness requires true \(X\).}

\hbrule{rule:p18-es-im-warrant}{Im-rule warrant (source schema)}
  {Q\qquad X\qquad\operatorname{Im}_L(Q,X,R)}{\text{raise }R}
  {Actual two-clause relation with proper nonempty witnesses; requires sound
   Q and true X for soundness. Earlier question need not have been queried.}

\hbrule{rule:p18-es-meeting}{Figure 5: meeting (imperative)}
  {Q=?\{\neg A,A\land B,A\land\neg B\}}
  {\text{ask }?A;\quad\begin{array}{l}\neg A\mapsto\neg A\\
   A\mapsto\text{ask }?B;\quad B\mapsto A\land B,\ \neg B\mapsto A\land\neg B\end{array}}
  {Empty basis. Second query uses original Q and accepted A.
   All three branches: \ref{annex:p18.es.meeting}.}

\hbrule{rule:p18-es-departure}{Figure 6: departure (imperative)}
  {Q=?\{A,B,C,D\};\quad E\supset(A\lor B),\ \neg E\supset D,\ B\equiv F}
  {\text{ask }?E;\quad\begin{array}{l}\neg E\mapsto D\\
   E\mapsto\text{ask }?F;\quad F\mapsto B,\ \neg F\mapsto A\end{array}}
  {Full positive-E branch has unqueried \(?\{A,B\}\).
   No C-endpoint: incomplete but fair. \ref{annex:p18.es.destination}.}

\hbrule{rule:p18-es-predicate}{Figure 7: predicate (imperative)}
  {?Pc;\quad\forall x(Px\equiv(Rx\land Tx))}
  {\text{ask }?Rc;\quad\begin{array}{l}\neg Rc\mapsto\neg Pc\\
   Rc\mapsto\text{ask }?Tc;\quad Tc\mapsto Pc,\ \neg Tc\mapsto\neg Pc\end{array}}
  {Distinct predicates; closed instances; nonempty domain. Full scenario has
   \(?\pm|Rc,Tc|\). \ref{annex:p18.es.predicate}.}

\hbrule{rule:p18-es-negation}{Figure 8: negation (imperative)}
  {?\neg A}
  {\text{ask }?A;\quad A\mapsto\neg\neg A,\ \neg A\mapsto\neg A}
  {Boolean double negation. Both responses: \ref{annex:p18.es.search-procedures}.}

\hbrule{rule:p18-es-conjunction}{Figure 9: conjunction (imperative)}
  {?(A\land B)}
  {\text{ask }?A;\quad\begin{array}{l}\neg A\mapsto\neg(A\land B)\\
   A\mapsto\text{ask }?B;\quad B\mapsto A\land B,\ \neg B\mapsto\neg(A\land B)\end{array}}
  {Boolean truth; full scenario retains binary question; no pure shortcut to
   \(?A\). \ref{annex:p18.es.search-procedures}.}

\hbrule{rule:p18-es-implication}{Figure 10: implication (imperative)}
  {?(A\supset B)}
  {\text{ask }?A;\quad\begin{array}{l}\neg A\mapsto A\supset B\\
   A\mapsto\text{ask }?B;\quad B\mapsto A\supset B,\ \neg B\mapsto\neg(A\supset B)\end{array}}
  {Boolean material conditional. Full scenario retains binary question.
   Negative-A branch is positive. \ref{annex:p18.es.search-procedures}.}

\hbrule{rule:p18-es-disjunction}{Figure 11: disjunction (imperative)}
  {?(A\lor B)}
  {\text{ask }?A;\quad\begin{array}{l}A\mapsto A\lor B\\
   \neg A\mapsto\text{ask }?B;\quad B\mapsto A\lor B,\ \neg B\mapsto\neg(A\lor B)\end{array}}
  {Boolean truth; full scenario retains binary question. All three routes:
   \ref{annex:p18.es.search-procedures}.}

\hbrule{rule:p18-es-equivalence}{Figure 12: equivalence (imperative)}
  {?(A\equiv B)}
  {\text{ask }?A;\ \text{then ask }?B;\quad\begin{array}{ll}
   A,B\mapsto A\equiv B&A,\neg B\mapsto\neg(A\equiv B)\\
   \neg A,B\mapsto\neg(A\equiv B)&\neg A,\neg B\mapsto A\equiv B\end{array}}
  {All four Boolean responses; repeated query labels are different prefix
   occurrences. Binary premise retained in full scenario.
   \ref{annex:p18.es.search-procedures}.}
